\documentclass[10pt,a4paper]{amsart}
\usepackage[margin=19mm]{geometry}
\usepackage{amsmath,amssymb,mathtools}
\usepackage[scr=rsfs,cal=euler]{mathalfa}
\usepackage{xfrac}
\usepackage[unicode,hidelinks,bookmarks=false]{hyperref}
\newcommand{\ie}{i.e.\ }
\newcommand{\cf}{cf.\ }
\newcommand{\resp}{respectively\ }
\newcommand{\Subsection}{\subsection}
\providecommand{\nobreakdash}{\nobreak\hbox{-}}
\setlength{\emergencystretch}{2em}
\multlinegap0pt
\usepackage{tikz}
\usepackage{amssymb}
\usepackage{float}
\usepackage{xcolor}
\usepackage[shortlabels]{enumitem}
\usepackage{enumerate}
\newtheorem{lemma}{Lemma}
\def\Exc{\mathsf{Exc}}
\def\Excp{\mathsf{Excp}}
\def\Nexc{\mathsf{Nexc}}
\def\den{\mathsf{den}}
\def\imv{\mathsf{imv}}
\def\inv{\mathsf{inv}}
\title{Further refinements of Euler-Mahonian statistics for multipermutations}
\author{Kaimei Huang, Yongzhou Wen, Sherry H.F. Yan}
\date{Source: arXiv:2601.20201v1 (CC BY 4.0)}
\begin{document}
\maketitle

\noindent\emph{Source and licence.} Further refinements of Euler-Mahonian statistics for multipermutations. Version arXiv:2601.20201v1 (CC BY 4.0), distributed by arXiv under the Creative Commons Attribution 4.0 International licence, http://creativecommons.org/licenses/by/4.0/, as recorded in the arXiv licence metadata for this exact version and shown on the abstract page https://arxiv.org/abs/2601.20201v1. Full licence notice: \texttt{licenses/arxiv-2601.20201-CC-BY-4.0.txt}.\par\smallskip

\noindent\textit{Editorial scope.} Counted unit: the article's lemma lem-phi-den1 (source0.tex lines 1090--1105). The article's complete original proof of it is delivered with it (source0.tex lines 1106--1238, one \texttt{proof} environment of 133 lines). No mathematics is written, repaired or re-derived: the statement and the proof are the article's own lines, in the article's own notation, and no retained line is substituted or re-flowed.  No further source-local context is needed: every cross-reference of every retained line resolves inside the two delivered spans.  Nothing was omitted from any retained span, so the package carries no omission record.  No retained line contains a citation, so the wrapper carries no bibliography and no bibliography entry.  No retained line contains a figure environment, an image-inclusion command or drawing code, so no image is delivered and the package declares no asset dependency.  Editorial formatting only, in addition: an AMS \texttt{amsart} preamble that restates the article's own declarations and macros used by the retained lines, namely the environments \texttt{lemma} on separate counters, together with the standard packages those lines need.  The retained environments print in this document as Lemma 1 in order of appearance, whereas the article prints the same items with the numbering of its own full document.  The complete native source of the article as distributed in the arXiv e-print for this exact version is bundled unchanged as \texttt{sources/193510/source0.tex}.
\begin{lemma}\label{lem-phi-den1}
Let $(w, c)\in  \mathfrak{S}_{N'}\times  \{0,1, \ldots, m'\}$ with $w=\alpha_1\alpha_2\ldots \alpha_{m'+a-1}$ and  $\bar{w}=x_1x_2$ $\ldots x_{m'+a-1}$. Assume that $w$ has exactly $s$ $g$-gap $h$-level excedance places.
Choose the integer $y$ such that    the space before  the letter  $\alpha_y$ is labeled by $c$  under the   $g\den_h$-labeling  of $w$ when $c>0$,  and let $y=m'+1$ when $c=0$.   Then we have 
\begin{equation}\label{eq-Excp_r}
	g\Excp_h(\phi_{n,g,h}(w, c))=\left\{ \begin{array}{ll}
		g\Excp_h(w)&\, \mathrm{if}\,\,   0\leq c\leq s+\gamma_g  \\
		g\Excp_h(w)\cup \{y\}&\, \mathrm{otherwise}.
	\end{array}
	\right.
\end{equation}
Moreover, 
	if  $\alpha_j\neq n$ for all $j<y$, then we have 
	\begin{equation}\label{eq-den_r}
g\den_h(\phi_{n,g,h}(w, c))=g\den_h(w)+c.
	\end{equation}
	\end{lemma}

\begin{proof}
We shall proceed the proofs of  (\ref{eq-Excp_r})	and (\ref{eq-den_r}) by considering the following three cases.

\noindent{\bf Case 1:} $c=0$. 
In this case, the multipermutation $u=\phi_{n,g,h}(w, 0)$ is obtained from $w$ by inserting an $n$ immediately after $\alpha_{m'}$. It is straightforward to see that we have $ 
g\Excp_h(u)=g\Excp_h(w).
$ 	If $\alpha_j\neq n$ for all $j\leq m'$, then all the $n's$ occurs to the right of $\alpha_{m'}$. This implies that each letter lying to the right of $\alpha_{m'}$ is equal to $n$. Then inserting  an $n$ immediately after $\alpha_{m'}$ does not affect  $g\den_{h}(w)$, implying that 
$g\den_h(u)=g\den_h(w)$ as desired. 


\noindent{\bf Case 2:}  $c>0$ and $y\notin g\Excp_h(w)$.\\
Recall that $u=\phi_{n,g,h}(w,c)$ is constructed  as follows.
	First, find all the  non-$g$-gap-$h$-level excedance letters of $w$ that occur    to  the right  of $\alpha_{y}$, say $\alpha_{j_1}, \alpha_{j_2}, \ldots, \alpha_{j_b}$ with $ j_1<j_2<\cdots<j_b$.  Then  replace $\alpha_{y}$ with an  $n$,   and replace  $\alpha_{j_{z+1}}$ with $\alpha_{j_{z}}$ for all $0\leq z\leq b $ with the convention that $j_{b+1}=m'+a$ and $j_0=y$.
	
	It is easily seen that  the procedure  from $w$ to $u$ keeps  all the $g$-gap $h$-level excedance letters of $w$  in   place.  Moreover,   this  procedure  also transforms $y$ into   a new  $g$-gap $h$-level excedance place when $y\leq m'-\gamma_g$ since for all $y\leq m'-\gamma_g$, we have $n\geq x_{y}+g$. 
	By the rules specified in the $g\den_h$-labeling of $w$, we have $c\leq \gamma_g$ when $y>m'-\gamma_g$,  and   $c\geq s+\gamma_g+1$ when $y\leq m'-\gamma_g$ and $y\notin g\Excp_h(w)$.  Therefore, we deduce that
	 $$
	 	g\Excp_h(u)=\left\{ \begin{array}{ll}
	 	g\Excp_h(w)&\, \mathrm{if}\,\,   1\leq c\leq s+\gamma_g  \\
	 	g\Excp_h(w)\cup \{y\}&\, \mathrm{otherwise},
	 \end{array}
	 \right.
	 $$
	as stated in (\ref{eq-Excp_r}). 
	       
Now we proceed to verify  (\ref{eq-den_r}).  We have two subcases. 

\noindent{\bf Subcase 2.1:} $y\leq  m'-\gamma_g$.\\
 Assume that there are exactly  $z$ $g$-gap $h$-level excedance letters that are located to the right of $\alpha_y$. 
By the rules   specified in the $g\den_h$-labeling of $w$,  one can easily check that $c=y+z+\gamma_g>s+\gamma_g$. In order to verify (\ref{eq-den_r}), it suffices to show that $g\den_h(u)=g\den_h(w)+y+z+\gamma_g$. 
Recall that 
$$
g\den_h(w)=\sum\limits_{i\in g\Excp_h(w)} (i+B^g_i(w)) +{\imv}(g\Exc_h(w))+{\inv}(g\Nexc_h (w)),
$$
	where $B^{g}_i(w)=|\{j \mid \alpha_i -g< x_j < \alpha_i\}|$. 
By (\ref{eq-Excp_r}), we have $g\Excp_h(u)=g\Excp_h(w)\cup \{y\}$.  By the construction of $u$, the procedure from $w$ to $u$ preserves all the non-$g$-gap-$h$-level excedance letters of $w$ as well as their relative order.  This implies that $g\Nexc_h(w)=g\Nexc_h(u)$.    Again by the construction of $u$, the procedure from $w$ to $u$ keeps all the $g$-gap $h$-level excedance letters of $w$ in place. 
Therefore, the procedure from $w$ to $u$  increases the sum 
$ 
\sum\limits_{i\in g\Excp_h(w)} (i+B^g_i(w))
$ by $y+B^g_y(u)=y+\gamma_g$ and does not affect ${\inv}(g\Nexc_h (w))$. Clearly, replacing $\alpha_y$ by an $n$ increases ${\imv}(g\Exc_h(w))$  by $z$ since there does not exist any $n's$ occurring to the left of $\alpha_y$ and there are exactly $z$ $g$-gap $h$-level excedance letters that are located to the right of $\alpha_y$ in $w$. Thus, we derive that $g\den_h(u)=g\den_h(w)+y+z+\gamma_g=g\den_h(w)+c$ as desired. 


\noindent{\bf Subcase 2.2:} $y>  m'-\gamma_g$.\\
By the rules   specified in the $g\den_h$-labeling of $w$,  one can easily check that $c=m'+1-y\leq \gamma_g$. In order to verify (\ref{eq-den_r}), it suffices to show that $g\den_h(u)=g\den_h(w)+m'+1-y$. 
Recall that 
$$
g\den_h(w)=\sum\limits_{i\in g\Excp_h(w)} (i+B^g_i(w)) +{\imv}(g\Exc_h(w))+{\inv}(g\Nexc_h (w)).
$$
   By the construction of $u$, the procedure from $w$ to $u$ keeps all the $g$-gap $h$-level excedance letters in place.  By (\ref{eq-Excp_r}), we have $g\Excp_h(u)=g\Excp_h(w)$.  It follows that   the procedure from $w$ to $u$ does not affect the sum $\sum\limits_{i\in g\Excp_h(w)} (i+B^g_i(w))$ and $\imv(g\Exc_h(w))$.  
   Assume that there are exactly $z$ letters that are smaller than $n$  and   occur weakly to the  right of $\alpha_y$ in $w$. Then   replacing $\alpha_y$ by an $n$  would increase  ${\inv}(g\Nexc_h(w))$  by $z$. 
   Hence we have $g\den_h(u)=g\den_h(w)+z$. Recall that there are no $n's$ that are located  to the left of $\alpha_y$ in $u$. This implies that $y=m'+1-z$. Therefore, we derive that
   $$
   g\den_h(u)=g\den_h(w)+z=g\den_h(w)+m'+1-y=g\den_h(w)+c
   $$
   as desired.  
   
   
\noindent{\bf Case 3:}  $c>0$ and $y\in g\Excp_h(w)$.\\
Recall that   $u=\beta_1\beta_2\ldots\beta_{m'+a}=\phi_{n,g,h}(w,c)$ is obtained  from $w$ as follows.
	\begin{itemize}
	\item  Find all the $g$-gap $h$-level excedance letters   of $w$  occurring      weakly to the right of   $\alpha_y$,  say
	$\alpha_{i_1}, \alpha_{i_2}, \ldots, \alpha_{i_a}$ with
	$ y=i_1<i_2<\cdots <i_a$. Find the smallest  integer $k$ satisfying $\alpha_{i_k}< x_{i_{k+1}}+g$   with the convention that $x_{i_{a+1}}=n$.
	\item    
	Choose the smallest integer  $p$ such that $x_p=\alpha_{i_k}-g+1$.	Find all the  non-$g$-gap-$h$-level excedance letters of $w$ that occur   weakly   to  the right  of $\alpha_{p}$, say $\alpha_{j_1}, \alpha_{j_2}, \ldots, \alpha_{j_b}$ with $ j_1<j_2<\cdots<j_b$.    
	\item   Replace  $\alpha_{i_1}=\alpha_y$ with an  $n$ and replace   $\alpha_{i_{z}}$  with $\alpha_{i_{z-1}}$  for all $1< z\leq k$. 
	
	\item  Replace  $\alpha_{j_1}$ with $\alpha_{i_k}$,  and  replace  $\alpha_{j_{z+1}}$ with $\alpha_{j_{z}}$ for all $1\leq z\leq b $  with the convention that $j_{b+1}=m'+a$.   
\end{itemize}

 Assume that there are exactly  $z$ $g$-gap  $h$-level excedance letters that are located to the right of $\alpha_y$. 
By the rules specified in the $g\den_h$-labeling of $w$, we have $c=z+1+\gamma_g\leq \gamma_g+s$.  In order to verify (\ref{eq-Excp_r}), it suffices to show that  $g\Excp_h(w)=g\Excp_h(u)$. 
According to the choices   of $k$ and $p$, we have $\alpha_{i_z}\geq  x_{i_{z+1}}+g$ for all $z<k$ and $\alpha_{i_k}< x_{j_1}+g$. This guarantees  that   the procedure from $w$ to $u$ keeps  all  the $g$-gap $h$-level excedance places, yielding that $g\Excp_h(w)=g\Excp_h(u)$ as desired.    
 

Recall that $c=z+1+\gamma_g$. 
  In order to verify (\ref{eq-den_r}), it suffices to show that $g\den_h(u)=g\den_h(w)+z+1+\gamma_g$. 
Recall that 
$$
g\den_h(w)=\sum\limits_{i\in g\Excp_h(w)}(i+B^g_i(w)) +{\imv}(g\Exc_h(w))+{\inv}(g\Nexc_h (w)).
$$
Since $g\Excp_h(u)=g\Excp_h(w)$,  by the construction of $u$, the procedure from $w$ to $u$ transforms $\beta_{y}=n$ into a new $g$-gap $h$-level excedance letter, transforms $\alpha_{i_k}$ into a non-$g$-gap-$h$-level excedance letter, and preserves the other $g$-gap $h$-level excedance letters. Recall that 
$\beta_{i_j}=\alpha_{i_{j-1}}$ for all $1<j\leq k$. This implies that $B^g_{i_j}(u)=B^g_{i_{j-1}}(w)$ for all $1<j\leq k$. 
 Hence, we deduce that 
\begin{equation}\label{eq-sum}
	\begin{array}{lll}
\sum\limits_{i\in g\Excp_h(u)} (i+B^g_i(u))&=&B^g_{y}(u)-B^g_{i_k}(w)+\sum\limits_{i\in g\Excp_h(w)} (i+B^g_i(w))\\
&=&\gamma_g-B^g_{i_k}(w)+\sum\limits_{i\in g\Excp_h(w)} (i+B^g_i(w)),
\end{array}
\end{equation} 
  where  the second equality follows from the fact that  $B^g_y(u)=\gamma_g$.  
  Recall that there are no occurrences of $n$ to the left of $\alpha_y$. 
By the  equality $g\Excp_h(u)=g\Excp_h(w)$,  one can easily check that the procedure from $w$ to $u$ preserves the number of $g$-gap $h$-level excedance letters that are located to the right of $\alpha_y$. This implies that replacing $\alpha_{y}$ with  an $n$ would increase  ${\imv}(g\Exc_h(w))$ by $z$. 

Assume that there are exactly $t$ $g$-gap $h$-level excedance letters that are located to the left of $\alpha_{i_{k}}$ and  are greater than or equal to $\alpha_{i_k}$. Again by the choice  of $i_k$, all the $g$-gap $h$-level excedance letters located to the right of $\alpha_{i_k}$ are greater  than $\alpha_{i_k}$.  Recall that in the  the procedure from $w$ to $u$, we  transforms the letter  $\alpha_{i_k}$ into a non-$g$-gap-$h$-level excedance letter.    This procedure would decrease ${\imv}(g\Exc_h(w))$ by $t$.   Hence, we deduce that \begin{equation}\label{eq-imv}
	{\imv}(g\Exc_h(u))={\imv}(g\Exc_h(w))+z-t.
\end{equation}
In view of  (\ref{eq-sum}) and (\ref{eq-imv}), in order to prove $g\den_h(u)=g\den_h(w)+z+1+\gamma_g$, it remains to show that 
\begin{equation}\label{eq-inv}
	{\inv}(g\Nexc_h(u))={\inv}(g\Nexc_h(w))+t+1+B^g_{i_k}(w).
\end{equation} 
 Note that the procedure from $w$ to $u$  preserves    all the non-$g$-gap-$h$-level excedance letters of $w$ as well as their relative order,  and transforms  $\alpha_{i_k}$ into a non-$g$-gap-$h$-level excedance letter.  In order to prove (\ref{eq-inv}),  it suffices to  verify the following facts.

\noindent{\bf Fact 1:} All the non-$g$-gap-$h$-level excedance letters lying  to the left of $\beta_{j_{1}}$ are smaller than    $\beta_{j_{1}}=\alpha_{i_k}$ in $u$.
 
  \noindent{\bf Fact 2:} There are exactly $t+1+B^g_{i_k}(w)$ non-$g$-gap-$h$-level excedance letters that are located to the right of $\beta_{j_{1}}$ and   are smaller than  $\beta_{j_{1}}=\alpha_{i_k}$ in $u$.
  
 
  Since $i_k\in g\Excp_{h}(w)$,  it follows that all the non-$g$-gap-$h$-level excedance letters occurring  to the left of $\alpha_{i_{k}}$  are smaller than $\alpha_{i_k}$.  In order to prove Fact $1$, it suffices  to verify the following two claims. 
  
  \noindent{\bf Claim 1:} For all $i_k<j<p$, we have $j\notin g\Excp_h(w)$.\\
   If not,  choose $q$ be the smallest  integer such that $i_k<q<p$ and 
  $q\in g\Excp_h(w)$.  Then we have $i_{k+1}=q$. According to the choice of $p$, we have $x_q\leq \alpha_{i_k}-g$. Then we have $\alpha_{i_k}\geq x_q+g=x_{i_{k+1}}+g$,  yielding a contradiction with the choice of $i_k$ .  Hence, we conclude that $j\notin g\Excp_{h}(w)$ for all $i_k<j<p$ as claimed.
  
  \noindent{\bf Claim 2:} For all  $i_k<j<p$, we have $\alpha_j< \alpha_{i_k}$. \\
  If not, assume that $\alpha_q\geq \alpha_{i_k}$ for some $i_k<q<p$.  This combined with Claim  $1$ yields that  $\alpha_{i_k}\leq \alpha_q\leq x_{q}+g-1$, implying  that $x_q\geq \alpha_{i_k}+g-1$. This contradicts     the choice of $p$, completing  the proof of Claim $2$. 
  
   By the construction of $u$, in order to prove  Fact 2, it suffices to show that there are exactly $t+1+B^g_{i_k}(w)$ letters that are located weakly to the right of  $\alpha_{j_{1}}$ and  are smaller than $\alpha_{i_k}$ in $w$. 
  The following claim follows directly from   the selection of $j_{1}$. \\
  \noindent{\bf  Claim 3:} For all $p\leq j<j_{1}$, we have $j\in g\Excp_{h}(w)$. \\
   Assume that 
   $$
   u=|\{j\mid x_p\leq \alpha_{j}<\alpha_{i_k}, j<p\} |  
   $$
   and 
   $$
   v=|\{j\mid x_p\leq \alpha_{j}<\alpha_{i_k}, j\geq j_{1}\} | .
   $$
   Recall that $B^g_{i_k}(w)=|\{j\mid \alpha_{i_k}-g+1\leq x_{j}<\alpha_{i_k} \}|=|\{j\mid x_p\leq \alpha_{j}<\alpha_{i_k} \}|$.
   Claim  3 ensures that $u+v=B^g_{i_k}(w)$.  
   Recall that there are exactly $t$ $g$-gap $h$-level excedance letters that are located to the left of $\alpha_{i_k}$ and   are greater than or equal to $\alpha_{i_k}$ in $w$.  This   combined with Claim 2 tells us that  there are exactly $t+1$ letters that are greater than or equal to $\alpha_{i_k}$ and   are located to the left of $\alpha_{p}$ in $w$.    So there are exactly $u+t+1$ letters that are located to the left of $\alpha_p$ and  are greater than or equal to $x_p$. This implies that there are exactly $u+t+1$ letters that are located weakly to the right of  $\alpha_{p}$ and  are smaller than $x_p=\alpha_{i_k}-g+1$. Hence, there are exactly $u+t+1+v=t+1+B^g_{i_k}(w)$ letters that are located weakly to the right of  $\alpha_{p}$ and  are smaller than $\alpha_{i_k}$. 
    By Claim 3, such $t+1+B^g_{i_k}(w)$   letters  occur weakly to the right of $\alpha_{j_{1}}$ in $w$, completing the proof of Fact 2. 
	\end{proof}

\end{document}
